% Reconstructed from the complete geometry proof audited before the outage.
% Root, B and D: complete proof and Conrad/Abbes--Saito source audits PASS.
\begin{proposition}[Proved: Explicit simple covers with uniformly bounded breaks]
\label{D:psl2-bounded-conductor-covers}
Let $k$ be algebraically closed of characteristic two and let $q=2^r\ge4$.
There are arbitrarily many jointly linearly disjoint
$S_q=\operatorname{PSL}_2(\mathbf F_q)$-covers of $\mathbf G_{m,k}$.
Each compactifies to a map between projective lines with exactly two
branch values. At one value the inertia is tame cyclic of order $q+1$;
at the other it is the Borel group of order $q(q-1)$, with lower groups
\[
 G_0=\mathbf F_q\rtimes\mathbf F_q^\times,
 \qquad G_1=\mathbf F_q,
 \qquad G_2=1.
\]
Thus its positive upper break is $1/(q-1)\le1/3$. The groups $S_q$ are
nonabelian simple and admit three explicit generators.
\end{proposition}
\begin{proof}
In characteristic two every projective linear transformation over
$\mathbf F_q$ has a unique determinant-one representative. Thus
$S_q=\operatorname{SL}_2(\mathbf F_q)=\operatorname{PGL}_2(\mathbf F_q)$,
of order $q(q^2-1)$. Choose a generator $\lambda$ of
$\mathbf F_q^\times$. The matrices
\[
 u(1)=\begin{pmatrix}1&1\\0&1\end{pmatrix},\qquad
 d(\lambda)=\begin{pmatrix}\lambda&0\\0&\lambda^{-1}\end{pmatrix},
 \qquad w=\begin{pmatrix}0&1\\1&0\end{pmatrix}
\]
generate $S_q$. Conjugating $u(1)$ by powers of $d(\lambda)$ gives
$u(\lambda^{2j})$; these parameters span $\mathbf F_q$ additively.
Conjugation by $w$ gives all lower elementary matrices, and elementary
row reduction generates the entire special linear group.

We include the simplicity argument to specify exactly what is used
\citep[Theorem 2.10 and Theorems 3.1--3.5, pp.~2--5]{C_ConradPSL}.
The action on $\mathbf P^1(\mathbf F_q)$ is faithful and doubly
transitive. Hence every nontrivial normal subgroup $N$ is transitive.
Let $U$ be the translation subgroup in the stabilizer $B$ of infinity;
it is abelian and normal in $B$, and its conjugates generate $S_q$.
For each $g$ choose $n\in N$ with $g\infty=n\infty$; then $n^{-1}g\in B$.
In $S_q/N$ every conjugate of the image of $U$ is consequently that
same abelian image. The quotient is abelian. On the other hand
\[
 [d(\lambda),u(b)]=u((\lambda^2-1)b)
\]
and $\lambda^2\ne1$, so the upper and lower elementary groups lie in
the commutator group. Thus $S_q$ is perfect, and $N=S_q$.

Put $A=x^q-x$. The rational function
\[
 J_q(x)=\frac{(x^{q^2}-x)^{q+1}}{(x^q-x)^{q^2+1}}
       =\frac{(A^{q-1}+1)^{q+1}}{A^{q(q-1)}}
\]
is invariant under $S_q$. Indeed for a Mobius transformation
$g(x)=(ax+b)/(cx+d)$ over $\mathbf F_q$,
\[
 g(x)^{q^a}-g(x)
 =\frac{\det(g)(x^{q^a}-x)}{(cx+d)^{q^a+1}}.
\]
The exponents cancel in the displayed quotient. The identity
$x^{q^2}-x=A(A^{q-1}+1)$ proves the second formula. Its numerator and
denominator are coprime, and the reduced degree is
$q(q-1)(q+1)=|S_q|$. Hence $k(x)^{S_q}=k(J_q)$; the quotient is
separable because it is the fixed field of a finite automorphism group.

The poles are exactly $\mathbf P^1(\mathbf F_q)$, and the zeros are
$\mathbf P^1(\mathbf F_{q^2})\setminus\mathbf P^1(\mathbf F_q)$.
At infinity the stabilizer is $B$, of order $q(q-1)$. At a quadratic
point its stabilizer is the nonsplit torus of order $q+1$: identifying
the two-dimensional vector space with $\mathbf F_{q^2}$ realizes it as
$\mathbf F_{q^2}^\times/\mathbf F_q^\times$. Every nonidentity Mobius
transformation has its fixed points in a field extension of degree at
most two. These are therefore the only points with nontrivial
stabilizer, and there are exactly two branch values.

Use $u=1/x$ at infinity. A translation $x\mapsto x+b$, $b\ne0$,
acts by $u\mapsto u/(1+bu)$; its difference from $u$ has order two.
A stabilizer element with nontrivial multiplier has difference of
order one. Hence $G_1$ consists exactly of translations and $G_2=1$.
Herbrand's formula gives upper break $1/(q-1)$. The residue field is
perfect, so this is also the logarithmic break; the corresponding
nonlogarithmic bound is $1+1/(q-1)\le4/3$, by the classical comparison
in Abbes--Saito Proposition 3.15.

For $c\in k^\times$ consider the quotient map $t=cJ_q(x)$, restricted
over $\mathbf G_m$. We show that distinct $c$ give inequivalent
extensions of $k(t)$. The normalizer of $S_q$ in
$\operatorname{PGL}_2(k)$ is exactly $S_q$. A normalizing transformation
preserves the fixed-point set of the elements of order two, which is
$\mathbf P^1(\mathbf F_q)$. Its images of $0,1,\infty$ determine it,
so it is defined over $\mathbf F_q$. It therefore fixes $J_q$.
An isomorphism of two of the extensions would induce such a normalizer
element and would force the two constants to be equal.

Finally distinct Galois extensions with the same nonabelian simple
group are jointly linearly disjoint. We prove this by induction.
A subdirect subgroup of $S_q^m$ which is not the full product imposes,
by Goursat's lemma and simplicity, a diagonal identification between
two factors. Such an identification would identify the corresponding
extensions, contrary to their inequivalence. This proves the assertion
for arbitrarily large finite families, and hence the stated supply of
jointly independent covers.
\end{proof}

\begin{remark}[Finite verification and scope]
The exact certificate accompanying
Proposition~\ref{D:bounded-ramification-generation-obstruction} checks
$q=4,8,16$: finite-field arithmetic, the three generators, both orbits,
stabilizer orders, lower ramification, the invariant identity and the
multiplicity inequalities. The formulas above prove the assertions
for every $q$. These covers supply a counterexample to a proposed
group-theoretic implication from bounded ramification and linear
cohomology; they are not a counterexample to the proper-boundary
questions D.1 or D.2.
\end{remark}
